% Zdroj: PDF priklady-podzim-2022.pdf, fyzická str. 2. Značka: společné okruhy. Zařazeno: M1.
\section{Posloupnosti}
\szzotazka{podzim 2022}{3}{Posloupnosti (limita, omezenost)}

\noindent\textbf{Zadání.}\par
Nechť $(a_n) = (a_1, a_2, \dots)$ je posloupnost reálných čísel.
\begin{enumerate}
  \item Definujte, co znamená, že $(a_n)$ je \emph{konvergentní} (má vlastní limitu).
  \item Definujte, co znamená, že $(a_n)$ je \emph{neomezená}.
  \item Dokažte, že když $(a_n)$ je neomezená, pak není konvergentní.
\end{enumerate}

\medskip
\noindent\textbf{Náčrt řešení.}\par
\begin{enumerate}
  \item Posloupnost $(a_n)$ konverguje (má vlastní limitu $L \in \mathbb{R}$) právě tehdy, když
    \[
      \forall \varepsilon > 0 \ \exists n_0 \ \forall n \geq n_0 : \ |a_n - L| < \varepsilon.
    \]
  \item Posloupnost $(a_n)$ je \emph{omezená} právě tehdy, když existuje $K > 0$ takové, že $|a_n| \leq K$ pro všechna $n$. \emph{Neomezená} znamená, že omezená není, tj.
    \[
      \forall K > 0 \ \exists n : \ |a_n| > K.
    \]
  \item Dokážeme obměnou (kontrapozicí): ukážeme, že každá konvergentní posloupnost je omezená; neomezená posloupnost tedy nemůže konvergovat.

    Nechť $a_n \to L$. Pro $\varepsilon = 1$ existuje $n_0$ takové, že pro všechna $n \geq n_0$ platí $|a_n - L| < 1$, a tedy
    \[
      |a_n| = |a_n - L + L| \leq |a_n - L| + |L| < |L| + 1.
    \]
    Mimo těchto členů zbývá jen konečně mnoho prvních členů $a_1, \dots, a_{n_0 - 1}$. Položíme
    \[
      K = \max\{\, |a_1|, \dots, |a_{n_0 - 1}|,\ |L| + 1 \,\}.
    \]
    Pak $|a_n| \leq K$ pro všechna $n$, takže posloupnost je omezená. Je-li tedy $(a_n)$ neomezená, nemůže být konvergentní. \qed
\end{enumerate}

\medskip
\noindent\textit{(V~PDF bez náčrtu řešení; náčrt doplněn k~výkladu okruhu.)}
